% appendices2.tex — Appendices 2.A and 2.B for Chapter 2

\subsection*{Appendix 2.A: Asymptotics with Fixed Regressors}

To prove results similar in conclusion to Proposition~2.5, we add to Assumption~2.1 the following.

\begin{assumption}[2.A.1]
\label{ass:2.A.1}
$\mathbf{X}$ is deterministic. As $n \to \infty$,
$\mathbf{S}_{\mathbf{xx}} \to \bm{\Sigma}_{\mathbf{xx}}$, a nonsingular matrix
(conventional convergence here, because $\mathbf{X}$ is nonrandom).
\end{assumption}

\begin{assumption}[2.A.2]
\label{ass:2.A.2}
$\{\varepsilon_i\}$ is i.i.d.\ with $\E(\varepsilon_i) = 0$ and $\Var(\varepsilon_i) = \sigma^2$.
\end{assumption}

\begin{proposition}[2.A.1 (asymptotics with fixed regressors)]
\label{prop:2.A.1}
Under Assumptions~2.1, \ref{ass:2.A.1}, and \ref{ass:2.A.2}, plus an assumption
called the Lindeberg condition on $\{\mathbf{g}_i\}$ (where
$\mathbf{g}_i \equiv \mathbf{x}_i \cdot \varepsilon_i$), the same conclusions as in Proposition~2.5 hold.
\end{proposition}

Rather than stating the Lindeberg condition, we explain why such condition is
needed. To prove the asymptotic normality of $\mathbf{b}$, we would take (2.3.8) and show
that $\sqrt{n}\,\bar{\mathbf{g}}$ converges to a normal distribution. The technical difficulty is that
the sequence $\{\mathbf{g}_i\}$ is a sequence of nonrandom vectors. For example,
\[
\Var(\mathbf{g}_i) = \E(\varepsilon_i^2\,\mathbf{x}_i\mathbf{x}_i')
= \E(\varepsilon_i^2)\,\mathbf{x}_i\mathbf{x}_i'
= \sigma^2\,\mathbf{x}_i\mathbf{x}_i'
\]
is not constant across observations. Consequently, neither the Lindeberg--Levy CLT
nor the Martingale Differences CLT is applicable. Fortunately, however, there
is a generalization of the Lindeberg--Levy CLT to nonstationary processes with
nonconstant variance, called the \textbf{Lindeberg--Feller CLT}, which places a technical
restriction on the tail distribution of $\mathbf{g}_i$ to prevent observations with large
$\mathbf{x}_i$ from dominating the distribution of $\sqrt{n}\,\bar{\mathbf{g}}$.


\subsection*{Appendix 2.B: Proof of Proposition 2.10}

This appendix provides a proof of Proposition~\ref{prop:2.10}. We start with the expression
for $\hat{\gamma}_j$ given in \eqref{eq:2.10.12}. The first part of the proof is to find an expression that is
asymptotically equivalent.
%
\begin{align*}
\sqrt{n}\,\hat{\gamma}_j
&= \sqrt{n}\,\tilde{\gamma}_j
- \frac{1}{n}\sum_{t=j+1}^{n}
(\mathbf{x}_{t-j}\cdot\varepsilon_t + \mathbf{x}_t\cdot\varepsilon_{t-j})'\,
\sqrt{n}\,(\mathbf{b} - \bm{\beta}) \\
&\quad + \sqrt{n}\,(\mathbf{b} - \bm{\beta})'\Bigl(\frac{1}{n}\sum_{t=j+1}^{n}
\mathbf{x}_t\,\mathbf{x}_{t-j}'\Bigr)(\mathbf{b} - \bm{\beta}) \\
&\underset{a}{\sim}
\sqrt{n}\,\tilde{\gamma}_j
- \frac{1}{n}\sum_{t=j+1}^{n}
(\mathbf{x}_{t-j}\cdot\varepsilon_t + \mathbf{x}_t\cdot\varepsilon_{t-j})'\,
\sqrt{n}\,(\mathbf{b} - \bm{\beta}) \\
&\qquad \text{(since the last term above vanishes)}
\end{align*}
%
\begin{align*}
&\underset{a}{\sim}
\sqrt{n}\,\tilde{\gamma}_j
- \E(\mathbf{x}_{t-j}\cdot\varepsilon_t + \mathbf{x}_t\cdot\varepsilon_{t-j})'\,
\sqrt{n}\,(\mathbf{b} - \bm{\beta}) \\
&\qquad \Bigl(\text{since }
\frac{1}{n}\sum_{t=j+1}^{n}
(\mathbf{x}_{t-j}\cdot\varepsilon_t + \mathbf{x}_t\cdot\varepsilon_{t-j})
\pto \E(\mathbf{x}_{t-j}\cdot\varepsilon_t + \mathbf{x}_t\cdot\varepsilon_{t-j})\Bigr) \\
&= \sqrt{n}\,\tilde{\gamma}_j - \bm{\mu}_j'\,\sqrt{n}\,(\mathbf{b} - \bm{\beta})
\quad \text{where } \bm{\mu}_j \equiv \E(\mathbf{x}_t\cdot\varepsilon_{t-j}) \\
&\qquad \text{(since } \E(\mathbf{x}_{t-j}\cdot\varepsilon_t) = \mathbf{0}
\text{ by \eqref{eq:2.10.15})} \\
&= \sqrt{n}\,\tilde{\gamma}_j
- \bm{\mu}_j'\,\sqrt{n}\,\mathbf{S}_{\mathbf{xx}}^{-1}
\frac{1}{n}\sum_{t=1}^{n}\mathbf{x}_t\cdot\varepsilon_t \\
&= \sqrt{n}\,\frac{1}{n}\sum_{t=j+1}^{n}\varepsilon_t\varepsilon_{t-j}
- \bm{\mu}_j'\,\sqrt{n}\,\mathbf{S}_{\mathbf{xx}}^{-1}
\frac{1}{n}\sum_{t=1}^{n}\mathbf{x}_t\cdot\varepsilon_t
\quad \text{(by \eqref{eq:2.10.7})} \\
&\underset{a}{\sim}
\sqrt{n}\,\frac{1}{n}\sum_{t=1}^{n}\varepsilon_t\varepsilon_{t-j}
- \bm{\mu}_j'\,\sqrt{n}\,\bm{\Sigma}_{\mathbf{xx}}^{-1}
\frac{1}{n}\sum_{t=1}^{n}\mathbf{x}_t\cdot\varepsilon_t
\quad (\text{since } \mathbf{S}_{\mathbf{xx}} \pto \bm{\Sigma}_{\mathbf{xx}}) \\
&\qquad \Bigl(\text{since }
\sqrt{n}\,\frac{1}{n}\sum_{t=1}^{n}\varepsilon_t\varepsilon_{t-j}
- \sqrt{n}\,\frac{1}{n}\sum_{t=j+1}^{n}\varepsilon_t\varepsilon_{t-j} \\
&\qquad\qquad = \frac{1}{\sqrt{n}}\sum_{t=1}^{j}\varepsilon_t\varepsilon_{t-j}
\pto 0 \text{ for each } j\Bigr) \\
&\underset{a}{\sim}
\mathbf{c}_j'\sqrt{n}\,\bar{\mathbf{g}}_j,
\end{align*}
%
where
%
\begin{equation}
\mathbf{c}_j \equiv \begin{bmatrix} 1 \\ -\bm{\Sigma}_{\mathbf{xx}}^{-1}\bm{\mu}_j
\end{bmatrix}_{((K+1)\times 1)}, \quad
\bar{\mathbf{g}}_j \equiv \frac{1}{n}\sum_{t=1}^{n}\underset{((K+1)\times 1)}{\mathbf{g}_{jt}}, \quad
\mathbf{g}_{jt} \equiv \begin{bmatrix}
\varepsilon_{t-j}\varepsilon_t \\ \mathbf{x}_t\cdot\varepsilon_t
\end{bmatrix}_{((K+1)\times 1)}.
\label{eq:2.B.1}
\end{equation}
%
Thus, letting
\[
\hat{\bm{\gamma}} \equiv \begin{bmatrix}
\hat{\gamma}_1 \\ \vdots \\ \hat{\gamma}_p
\end{bmatrix}_{(p\times 1)},
\]
we have proved that
\[
\sqrt{n}\,\hat{\bm{\gamma}} \underset{a}{\sim}
\mathbf{C}'\sqrt{n}\,\bar{\mathbf{g}}
\]
where
\[
\underset{(p(K+1)\times p)}{\mathbf{C}}
\equiv \begin{bmatrix}
\mathbf{c}_1 & & \\ & \ddots & \\ & & \mathbf{c}_p
\end{bmatrix}, \quad
\underset{(p(K+1)\times 1)}{\bar{\mathbf{g}}}
\equiv \frac{1}{n}\sum_{t=1}^{n}\mathbf{g}_t, \quad
\underset{(p(K+1)\times 1)}{\mathbf{g}_t}
\equiv \begin{bmatrix}
\mathbf{g}_{1t} \\ \vdots \\ \mathbf{g}_{pt}
\end{bmatrix}.
\]

The second part of the proof is to show that $\mathbf{g}_t$ is a martingale difference
sequence, namely, that
\[
\E(\mathbf{g}_{jt} \mid \mathbf{g}_{t-1}, \mathbf{g}_{t-2}, \ldots\,) = \mathbf{0}
\quad (\text{for } j = 1, 2, \ldots, p).
\]
%
But this easily follows from the Law of Total Expectations, the fact that
$(\mathbf{g}_{t-1}, \mathbf{g}_{t-2}, \ldots)$ has less information than
$(\mathbf{x}_t, \mathbf{x}_{t-1}, \ldots, \varepsilon_{t-1}, \varepsilon_{t-2}, \ldots)$,
and \eqref{eq:2.10.15}.

Clearly, $\mathbf{g}_t$ is ergodic stationary. Therefore, by the ergodic stationary Martingale Differences CLT, $\sqrt{n}\,\bar{\mathbf{g}} \dto N(\mathbf{0}, \E(\mathbf{g}_t\mathbf{g}_t'))$.
The next step is to calculate $\E(\mathbf{g}_t\mathbf{g}_t')$. Its $(j, k)$ block is
%
\[
\underset{((K+1)\times(K+1))}{\E(\mathbf{g}_{jt}\mathbf{g}_{kt}')}
= \begin{bmatrix}
\E(\varepsilon_t^2\varepsilon_{t-j}\varepsilon_{t-k})
  & \E(\mathbf{x}_t'\varepsilon_t^2\varepsilon_{t-j}) \\
\E(\mathbf{x}_t\varepsilon_t^2\varepsilon_{t-k})
  & \E(\varepsilon_t^2\,\mathbf{x}_t\mathbf{x}_t')
\end{bmatrix}.
\]
%
By using the Law of Total Expectations and \eqref{eq:2.10.16}, it is easy to show that
%
\begin{equation}
\underset{((K+1)\times(K+1))}{\E(\mathbf{g}_{jt}\mathbf{g}_{kt}')}
= \begin{bmatrix}
\sigma^4\delta_{jk} & \sigma^2\bm{\mu}_k' \\
\sigma^2\bm{\mu}_j & \sigma^2\bm{\Sigma}_{\mathbf{xx}}
\end{bmatrix},
\label{eq:2.B.2}
\end{equation}
%
where $\delta_{jk}$ is ``Kronecker's delta,'' which is 1 if $j = k$ and 0 otherwise.

Since, as shown above,
$\sqrt{n}\,\hat{\bm{\gamma}} \underset{a}{\sim} \mathbf{C}'\sqrt{n}\,\bar{\mathbf{g}}$
and $\sqrt{n}\,\bar{\mathbf{g}} \dto N(\mathbf{0}, \E(\mathbf{g}_t\mathbf{g}_t'))$, we have
\[
\avar(\hat{\bm{\gamma}})
\;(\equiv \text{variance of limiting distribution of }
\sqrt{n}\,\hat{\bm{\gamma}})
= \mathbf{C}'\,\E(\mathbf{g}_t\mathbf{g}_t')\,\mathbf{C}.
\]
Its $(j, k)$ element is
\[
(j, k) \text{ element of } \avar(\hat{\bm{\gamma}})
= \mathbf{c}_j'\,\E(\mathbf{g}_{jt}\mathbf{g}_{kt}')\,\mathbf{c}_k.
\]

Substituting \eqref{eq:2.B.2} into this and using the definition of $\mathbf{c}_j$ in \eqref{eq:2.B.1}, we obtain,
after some simple calculations,
\[
(j, k) \text{ element of } \avar(\hat{\bm{\gamma}})
= \sigma^4 \cdot \bigl[\delta_{jk}
- \bm{\mu}_j'\,\bm{\Sigma}_{\mathbf{xx}}^{-1}\,\bm{\mu}_k/\sigma^2\bigr].
\]
So
\[
\sqrt{n}\,\hat{\bm{\gamma}} \dto
N(\mathbf{0},\;\sigma^4\mathbf{I}_p - \sigma^4\mathbf{\Phi}),
\]
where $\mathbf{\Phi}$ is defined in Proposition~\ref{prop:2.10}. Since
$\sqrt{n}\,\hat{\bm{\rho}} \underset{a}{\sim} \sqrt{n}\,\hat{\bm{\gamma}}/\sigma^2$,
the limiting distribution of $\sqrt{n}\,\hat{\bm{\rho}}$ is the same as that of
$\sqrt{n}\,\hat{\bm{\gamma}}/\sigma^2$, which is
$N(\mathbf{0}, \mathbf{I}_p - \mathbf{\Phi})$. \qed
